\section{Introduction}
\label{sec:intro}

A latent quantity is often observed only through a coarsening: a
device that maps the quantity of interest to a coarser report and
discards the rest. A failed component is reported as a candidate set
of suspects. A transcript count is reported after a dropout that
sometimes records a zero. A spot of tissue is reported as a mixture of
several cell types. A sensitive statistic is reported only after
calibrated noise is added. A label is reported as a vote of several
imperfect heuristics. A patient's true disease state is reported only
as a billing code. These are different fields with different
vocabularies, but the statistical object is the same: inference for a
latent parameter from data that have been coarsened by a mechanism the
analyst may or may not control.

\subsection{One idea, seven fields}

The masked-data likelihood framework of \citet{towell2026masked} gives
the common skeleton. A latent cause is observed as a candidate set;
three coarsening-at-random conditions, named C1, C2, C3
\citep{heitjan1991ignorability,gill1997coarsening,little2002statistical},
say when the coarsening mechanism factors out of the likelihood, so
the analyst may maximize a face-value likelihood without modeling the
mechanism; and an augmented-candidate-set rank condition says when the
latent parameter is identifiable from the coarsened data. That paper
treats series-system reliability. Six sibling papers carry the same
skeleton into other domains: scRNA-seq dropout
\citep{towell2026scrnacoarsening}, spatial-transcriptomics
deconvolution \citep{towell2026spatialcoarsening}, differential
privacy \citep{towell2026dpcoarsening}, programmatic weak supervision
\citep{towell2026weaksupcoarsening}, electronic-health-record
phenotyping \citep{towell2026phenotypecoarsening}, and multiple-instance
learning \citep{towell2026milcoarsening}. A robustness study
\citep{towell2026mdrelax} treats sensitivity to the symmetry condition
C2.

Each sibling, when it reaches its central structural result,
rediscovers the same theorem and gives it a local name. In scRNA-seq
it is \emph{cell-total consistency}: at the zero-inflated
negative-binomial maximum-likelihood estimate, the fitted mean
$(1-\hat\pi_j)\hat\mu_j$ equals the empirical mean $\bar X_j$. In
spatial transcriptomics it is the same identity at the spot level,
$\sum_k \hat P_{sk}\hat\mu_{j,k} = \bar X_{sj}$. In differential
privacy it is \emph{release consistency}: a mechanism-aware fit
reproduces the observed release in expectation. In weak supervision it
is \emph{agreement consistency}: the fitted label model reproduces the
empirical pairwise agreement rates of the labeling functions. In
phenotyping it is \emph{code-frequency consistency}: the fitted code
frequency $\hat\pi\,\widehat{\mathrm{sens}} +
(1-\hat\pi)(1-\widehat{\mathrm{spec}})$ equals the empirical code
frequency $\bar C$. In multiple instance learning it is
\emph{bag-prevalence consistency}: the fitted bag-positive probabilities
reproduce the observed bag-positive rates along the column space of the
bag composition matrix. The six identities have the same content and the
same practical moral, that marginal-fit agreement is not evidence of
unbiasedness, yet they are stated and proved six times.

A second result recurs the same way. Each sibling needs to know when
the latent parameter is recoverable, and answers with the
augmented-candidate-set rank condition of \citet{towell2026masked}.
And each sibling identifies the same remedy when identifiability
fails: a \emph{singleton} candidate set, a report that pins the latent
quantity to a point. The singleton wears a different costume in each
field, an ERCC spike-in, a single-cell-resolution probe, a gold label,
a non-private release, a chart-reviewed patient, a singleton bag, but it
is the same device.

\subsection{Contribution: state it once}

This paper states the recurring results once, at the generality that
explains the recurrence, and recovers each named version as a
specialization. The contributions are the following.

\begin{enumerate}[(i)]
\item A single general framework section (\cref{sec:framework}) that
  states C1, C2, C3 once for an abstract coarsening, defines the
  \emph{coarsening-sufficient statistic} that each domain instantiates,
  and gives the face-value likelihood that C1--C2--C3 license.
\item A general identifiability result (\cref{sec:identifiability}),
  the seam-free half of the synthesis. We state the
  augmented-candidate-set rank condition once and the
  singleton-restoration result once, and show that one $|c| = 1$ device
  restores the same column rank across all seven domains, with the
  per-domain singletons tabulated as instances of a single mechanism
  (\cref{tab:singletons}). This half carries no caveats and is not
  reducible to a textbook identity.
\item A general consistency theorem (\cref{thm:general-consistency},
  \cref{sec:consistency}), the more nuanced half, whose boundary we
  locate exactly. At an interior maximum of the face-value likelihood
  the fitted mean of the coarsening-sufficient statistic equals its
  empirical mean. We give the result in two regimes, a
  regular-exponential-family regime where the equality is exact in
  finite samples and a location-family regime where it is exact at a
  single report and a population first-moment identity otherwise, and
  we recover cell-total, spot-level, release, agreement,
  code-frequency, and bag-prevalence consistency as corollaries
  (\cref{cor:scrna,cor:spatial,cor:dp,cor:weaksup,cor:phenotype,cor:mil}),
  saying precisely which are exact, which are asymptotic, and which hold
  only under a richer parametrization.
\item A compact applications tour (\cref{sec:applications}) showing
  each domain is an instance, citing the sibling paper for the
  domain-specific development rather than re-deriving it.
\end{enumerate}

The claim is unification, not new mathematics. The six named
consistency theorems and the per-domain rank conditions already exist
in the siblings; the contribution is to show they are one theorem and
one condition, to state them at the right level, and to be precise
about where the unification is clean and where it is seamed. The count
is seven domains but six named consistency corollaries: reliability
\citep{towell2026masked} is the source of the shared framework and of
the singleton-and-rank apparatus, so it appears among the seven domains
and contributes a singleton device (\cref{tab:singletons}), but its
consistency content is the general identity itself rather than a
separately named corollary, which is why the consistency corollaries
number six. Two seams
matter and we flag them in advance. The weak-supervision corollary is
exact only for a parametrization in which the agreement statistics are
sufficient; the standard naive-Bayes label model satisfies the
identity only asymptotically. And differential privacy does not live
in the exponential-family branch at all: its release-consistency
identity is a first-moment statement for a continuous location family,
proved through the location-family score, and it is recovered from the
location-family branch of the general theorem rather than from the
exponential-family branch that serves the other five domains. Naming
these seams is part of the contribution, because a synthesis that
hides them would misrepresent what is shared.

One objection deserves a direct answer, because a reader will raise it
at once: in its exponential-family regime the consistency identity is
the textbook score equation $\sum_i \T(R_i) = n\,\E[\T]$ at the MLE, so
is the contribution merely that several papers reused an elementary
fact? It is not, for three reasons. First, the content of the identity
here is not the moment-matching itself but what the coarsening does to
it: the coarsening makes the \emph{observed} marginal an exponential
family while the estimation bias lives in the \emph{latent} parameter,
so the identity is a statement about what a fitted latent model
structurally cannot detect, not a computation. That is the diagnostic
moral every sibling issues, and it is the opposite of a triviality.
Second, the singleton-and-rank half of the paper is not elementary
moment-matching at all: the augmented-candidate-set rank condition and
the singleton-restoration result are a genuine identifiability device,
and showing that one $|c| = 1$ construction restores the same column
rank across seven domains is the harder and less attackable claim. Third,
the discovery that differential privacy does \emph{not} fit the
exponential-family box, and forces a structurally distinct
location-family regime with a precisely located boundary
(\cref{rem:loc-sketch}), is itself a result about the reach of the
principle rather than a restatement of a known one.

\subsection{Scope}

This is a methods-and-synthesis paper. It contains no new simulations
and no new data analysis; the empirical validation of each named
identity lives in the corresponding sibling, which we cite. The
mathematical content is a re-derivation at higher generality plus an
audit of which instances reduce cleanly. Where the general argument
needs a step that the siblings establish only in special cases, we
present the step as a proof sketch and say what remains to be filled
in for a fully general proof.
